% -------------------------------------------------------%
% パッケージのインポート

\usepackage{graphicx} % 図の挿入(includegraphics)
\usepackage{wrapfig}
\usepackage{booktabs,multirow,tabularx}
\usepackage{float} % [H]で厳密に位置を固定

\usepackage{bm}
\usepackage{amsmath} % 数式用
\usepackage{amssymb} % もっと数式記号(iintとか)
\usepackage{mathtools} % もっと数式(アクセント)
\usepackage{accents} % \undertildeで下付きチルダ
\usepackage{nccmath} % 数式左寄せ環境fleqn
\usepackage{empheq} % わからん

\usepackage{amsthm} % 定理，証明など

\usepackage{siunitx} % 単位書き方(si)
\sisetup{
  group-separator = {,},
  group-digits = integer,
  group-minimum-digits = 4,
  per-mode=symbol
}

\DeclarePairedDelimiter{\rbra}{\lparen}{\rparen}
\usepackage{physics2}
\usephysicsmodule{ab}
\usepackage[italic]{derivative}

\usepackage{url}

\usepackage{hyperref,theoremref,xcolor}
\definecolor{clink}{RGB}{0,60,110}%
\hypersetup{
  pdftitle={tohoku-civil-graduate-exam-maths},
  colorlinks=true, linkcolor=clink!90,
  bookmarksnumbered=true,
  bookmarksopen=true
}

\usepackage{enumitem} % enumerate
\setlist[enumerate,itemize]{
  itemsep=0pt,
  itemindent=4em,
  leftmargin=.5em,
  listparindent=1em,
  labelsep=.7em,
  itemindent=1.2em
}
% \usepackage{comment} % コメントアウト(comment環境)

\usepackage{framed}
\usepackage{ascmac}
\usepackage{fancybox}

% ページスタイルのため
\usepackage{fancyhdr}
\usepackage{lastpage}
\usepackage{color}

\usepackage{setspace}

% -------------------------------------------------------%

% ページ番号に西暦を入れる用のカウンター
\newcounter{yearcounter}

% yearcounterとpageを変更
% セクション始まりにこの2つを書けばページ番号を振り直せる
% \setcounter{yearcounter}{2020}
% \setcounter{page}{1}

% これは多分使わない
% yearcounterを1増加させる
\newcommand{\incyearcounter}{\stepcounter{yearcounter}}


% 1ならなにもしない、それ以外なら改ページする変数とコマンド
% cpornot
\newcounter{cpornot}
\setcounter{cpornot}{1}
\newcommand{\clearpageornot}{%
  \ifnum\value{cpornot}=1
    % skip
  \else
    \newpage
  \fi
}

% -------------------------------------------------------%
% 2:問題番号を表示しない 3:する
\setcounter{tocdepth}{2}
% \setcounter{tocdepth}{3}
